\section*{Limitations}

The current evidence base is intentionally narrow: the workshop-era experiments
cover a small asset set, limited scene diversity, and mostly visual or
feature-level metrics. The GRScenes extension adds real VLM grounding probes,
but the claim boundaries remain narrow. The clean visual-QA pool has 15 pairs,
below the configured 20-pair clean final gate, so it is not evidence that
conversion preserves VLM grounding in general. The 30-pair expanded stress set
is a frozen, target-centered material-shift set with 4 PASS and 26 WARN pairs;
it supports a controlled stress-test claim, not a broad GRScenes benchmark
distribution. Qwen2.5-VL also exposes unresolved coordinate semantics because
raw image-space scoring is stronger than the requested normalized-1000
interpretation. These results support claims about measurement protocols,
coordinate contracts, and model sensitivity, but they do not yet establish
broad downstream benefits for language, vision-language, or embodied-agent
models. The material-effect baseline is also claim-bounded: the current matrix
uses 30 GRScenes samples for four covered effect bins and two supplemental
fixtures for clearcoat and procedural texture. It supports selected qualitative
and failure-case interpretation, but it is not a material-effect error-rate
distribution. The added InternNav/DualVLN~\cite{Wei2026Ground} route is a
scoped official-route sanity check rather than a broad embodied benchmark: the
current official KuJiaLe \texttt{val\_unseen} evidence covers three local scenes
and 99 paired episodes, but it should not be promoted to all-InteriorNav,
R2R/MP3D, manipulation, or broad GRScenes robustness. The selected rollout
videos are qualitative evidence only. The official-scene load/render benchmark
also remains claim-bounded: original MDL and noMDL both complete all 18 required
fresh-process runs, but their warmed-ready intervals overlap, and no NVIDIA
official-scene baseline is available yet. The current pipeline also depends on
Isaac Sim and USD assets whose licenses, material graphs, and renderer settings
may differ across datasets, so any general guidelines should be validated on a
larger and more heterogeneous scene corpus.
